\def\stat{bor+bos}





\def\tit{АНАЛИЗ ФУНКЦИОНИРОВАНИЯ ИНФОРМАЦИОННО-АНАЛИТИЧЕСКИХ СИСТЕМ: 


МЕТОДИКА РАСЧЕТА ВРЕМЕННЫХ ЗАТРАТ И~ЧИСЛЕННОСТИ ПЕРСОНАЛА}





\def\titkol{Анализ функционирования ИАС: %информационно-аналитических систем: 


методика расчета временных затрат} % и~численности персонала}








\def\aut{А.\,В.~Борисов$^1$, А.\,В.~Босов$^2$, А.\,В.~Иванов$^3$, 


А.\,И.~Стефанович$^4$}





\def\autkol{А.\,В.~Борисов, А.\,В.~Босов, А.\,В.~Иванов, 


А.\,И.~Стефанович}





\titel{\tit}{\aut}{\autkol}{\titkol}





\index{Борисов А.\,В.}


\index{Босов А.\,В.}


\index{Иванов А.\,В.}


\index{Стефанович А.\,И.}


\index{Borisov A.\,V.}


\index{Bosov A.\,V.}


\index{Ivanov A.\,V.} 


\index{Stefanovich A.\,I.}








%{\renewcommand{\thefootnote}{\fnsymbol{footnote}}


%\footnotetext[1]


%{Работа выполнена при поддержке РФФИ (проект 15-07-02053).}}





\renewcommand{\thefootnote}{\arabic{footnote}}


\footnotetext[1]{Институт проблем информатики Федерального исследовательского центра 


<<Информатика и~управление>> 


Российской академии наук, \mbox{ABorisov@ipiran.ru}}


\footnotetext[2]{Институт проблем информатики Федерального исследовательского центра 


<<Информатика и~управление>>


Российской академии наук, \mbox{AVBosov@ipiran.ru}}


\footnotetext[3]{Институт проблем информатики Федерального исследовательского центра 


<<Информатика и~управление>>


Российской академии наук, \mbox{AIvanov@ipiran.ru}}


\footnotetext[4]{Институт проблем информатики Федерального исследовательского центра 


<<Информатика и~управление>>


Российской академии наук, \mbox{AStefanovich@frccsc.ru}}





 \label{st\stat}





                  \thispagestyle{headings}


                  


  \vspace*{-8pt}





        


        \Abst{Предложенный ранее подход к~оценке характеристик функционирования 


программного обеспечения на основе имитационного моделирования и~последующей 


статистической обработки результатов экспериментов развивается в~связи с~потребностями 


управления пользовательским персоналом ин\-фор\-ма\-ци\-он\-но-ана\-ли\-ти\-че\-ской  


сис\-те\-мы (ИАС). Для расчета чис\-лен\-ности персонала деятельность пользователей типизирована 


в~соответствии с~комплексом предметных задач и~функционалом привлекаемых программных 


средств. Представлены методические рекомендации по определению набора показателей 


временных затрат и~сбору необходимой статистической информации, предложены конкретные 


измеряемые показатели, приведены алгоритмы вычисления параметров трудоемкости выполнения 


операций и~производительности труда пользователей на основе собранных статистических 


данных. Разработанная методика оценки численности персонала ИАС проиллюстрирована 


простыми примерами.}


        


        \KW{информационно-аналитическая система; вероятностные характеристики; 


статистические методы; факторный анализ}





\DOI{10.14357\!/\!08696527180303} 








\vskip 10pt plus 9pt minus 6pt








      \thispagestyle{headings}


      


      \vspace*{-18pt} 


        


\section{Введение}





      В процессе создания и~сопровождения  


ин\-фор\-ма\-ци\-он\-но-те\-ле\-ком\-му\-ни\-ка\-ци\-он\-ной системы 


решается много задач, связанных с~определением характеристик ее 


функционирования. Если система создается с~целью поддержки принятия  


опе\-ра\-тив\-ных\!/\!так\-ти\-че\-ских\!/\!стра\-те\-ги\-че\-ских решений, то 


обязательными ее со\-став\-ля\-ющи\-ми являются~\cite{2-bor}:


      \begin{itemize}


\item локальные информационные массивы различной степени 


структуриро\-ван\-ности;


\item интерфейсы, обеспечивающие доступ к внешним информационным 


ресурсам и~сети Интернет;


\item средства управления пред\-мет\-но-ори\-ен\-ти\-ро\-ван\-ны\-ми базами 


данных (Data Warehouse);


\item средства интерактивной обработки данных (OLAP);


\item средства интеллектуального глубинного анализа данных (Data Mining);


\item специальные программные средства решения аналитических 


задач целевого назначения системы.


\end{itemize}





      Подобные информационно-телекоммуникационные системы 


называются ин\-фор\-ма\-ци\-он\-но-ана\-ли\-ти\-че\-ски\-ми~\cite{2-bor}. 


В~связи с~кардинальным 


усложнением задач, решаемых системами данного класса, проблема 


сопровождения ИАС значительно усложняется. Для ИАС в~полной мере 


значимы вопросы расчета надежности аппаратуры и~программного 


обеспечения, актуальны вопросы определения ве\-ро\-ят\-но\-ст\-но-вре\-мен\-н$\acute{\mbox{ы}}$х 


характеристик специального программного обеспечения (СПО). 





Для решения 


этих задач в~\cite{1-bor, 3-bor} изложены подходы к~соответствующим расчетам, 


апробированные на практике. Однако ключевую роль в~функционировании 


ИАС играет не аппаратное или программное обеспечение, а персонал. 


В~данном случае речь идет не о~представителях службы технической 


поддержки и~эксплуатации, а~о~пользователях ИАС, выполняющих целевые 


функции систем. В~отличие от необходимого числа сотрудников службы 


технической поддержки определить потребности ИАС 


в~\textit{пользовательском} персонале как величину, пропорциональную 


числу обслуживаемых уз\-лов\!/\!ре\-сур\-сов системы, не пред\-став\-ля\-ет\-ся 


возможным. Имеет место скорее обратная ситуация, когда численность 


пользователей ИАС определяет количество необходимых для их работы 


автоматизированных рабочих мест (АРМ).


      


      На эффективность и~качество работы пользовательского персонала 


влияет значительное число факторов, среди которых можно выделить:


      \begin{itemize}


\item квалификацию и~уровень подготовки пользователей;


\item заинтересованность и~ответственность пользователей;


\item степень однородности задач, решаемых каждым конкретным 


пользователем;


\item интенсивность работы в~ИАС (объем ежедневной нормы заданий);


\item текущие индивидуальные психофизические характеристики каждого 


пользователя;


\item необходимость совмещения пользователями выполнения функций 


в~ИАС и~других должностных обязанностей;





\pagebreak





\item потребительские качества инструментального пользовательского СПО 


в~составе ИАС;


\item характеристики и~состояние АРМ.


\end{itemize}





      Объективные же показатели, возможные к~определению в~рамках 


функционирования ИАС,~--- временн$\acute{\mbox{ы}}$е затраты на выполнение основных 


операций и~основанные на них показатели численности персонала~--- 


обсуждаются в~данной статье. Целью является разработка типовой методики 


определения численности пользовательского персонала, необходимого для 


обеспечения работы некоторой ИАС. Она состоит из следующих 


взаимосвязанных частей. 





Первая часть содержит классификацию операций 


в~ИАС, являющихся автоматизированными или ручными, т.\,е.\ 


нуждающимися в~большем или меньшем вмешательстве пользователей ИАС. 


Предлагаемая классификация может быть развита или скорректирована для 


каждой конкретной существующей или перспективной ИАС при наличии 


детальной информации о целях ее функционирования, архитектуре, 


применяемом про\-грам\-мно-ал\-го\-рит\-ми\-че\-ском, аппаратном 


и~коммуникационном обеспечении. Целевые функции ИАС достигаются 


путем выполнения того или иного набора представленных операций, 


а~иногда предполагают их итеративную или <<древовидную>> комбинацию.


      


      Вторая часть методики представляет простейшие математические 


модели для каждой предложенной операции ИАС. Она также содержит 


алгоритмы уточнения этих моделей на основе статистических данных 


функционирования ИАС с~использованием аппарата факторного анализа~[4]. 





Третья часть методики представляет формулы расчета численности 


пользовательского персонала в~зависимости от интенсивности заявок на 


выполнение тех или иных целевых функций ИАС.


      


      Статья организована следующим образом. Раздел~2 содержит 


классификацию операций в~ИАС, требующих участия пользователя. В~этом 


же разделе представлены:


      \begin{itemize}


\item математические модели, описывающие трудоемкость этих операций;


\item мероприятия, позволяющие собрать статистическую информацию об 


этих операциях в~процессе различного рода испытаний, опытной 


и~промышленной эксплуатации;


\item формулы обработки полученной статистической информации, 


позволяющие уточнить значения факторов в~математических моделях 


трудоемкости операций.


\end{itemize}





      Раздел~3 содержит методику расчета пользовательского персонала 


и~прос\-тей\-ший пример ее использования. В~разд.~4 представлены 


заключительные замечания и~возможные направления дальнейших 


исследований.





\section{Методика оценивания временных характеристик решения 


аналитических задач}





\subsection{Общие принципы оценивания временных характеристик}





      Настоящая методика базируется на следующих принципах.


      \begin{enumerate}[1.]


\item  Опыт создания самых разных ИАС позволяет обоснованно 


предполагать, что процесс решения любой целевой аналитической задачи 


представляет собой выполнение набора однотипных операций, которые 


могут быть классифицированы следующим образом:


\begin{itemize}


\item автоматизированные операции:


\begin{itemize}


\item ввод текстовых данных;


\item поиск\!/\!отбор результатов;


\item ввод и~обработка мультимедийных данных;


\item анализ\!/\!интерпретация\!/\!принятие решений по 


результатам выполнения поисков и~проведения расчетов;


\end{itemize}


\item автоматические операции:


\begin{itemize}


\item передача данных между узлами ИАС;


\item логистические операции при выполнении 


аналитических запросов;


\item аналитические расчеты;


\item формирование отчетов.


\end{itemize}


\end{itemize}


      В конкретной ИАС данный перечень может меняться или уточняться.


      


      Автоматизированные операции проводятся с~участием пользователей 


ИАС и~связаны с~действиями, требующими осмысленного вмешательства 


человека при вводе и~обработке разнородной информации, а также принятии 


решений. Автоматические операции в~полном объеме выполняются 


программным обеспечением ИАС без ка\-ко\-го-либо участия пользователей 


и~связаны с~приемом, обработкой и~хранением информации.


      


      \item Процесс выполнения функций назначения ИАС представляет 


собой решение некоторого набора целевых задач, который складывается из 


последовательности применения перечисленных выше операций. При этом 


для каждой задачи возможно существование нескольких сценариев ее 


выполнения, подразумевающих привлечение различного набора операций 


и~последовательности их применения. Более того, решение целевых задач 


часто представляет собой итеративный или древовидный процесс. Эти 


обстоятельства серьезно влияют на трудоемкость решения любой целевой 


задачи в~каждом конкретном случае.


      \item Временн$\acute{\mbox{ы}}$е затраты на выполнение автоматизированных 


операций зависят от следующих факторов, которые предлагается учитывать 


при расчете трудоемкости:


      \begin{itemize}


\item индивидуального уровня подготовки пользователя ИАС: опыта работы 


со средствами вычислительной техники вообще, конкретным СПО ИАС 


и~пр.;


\item качества имеющейся в~ИАС информации, т.\,е.\ ее содержательности, 


полноты и~непротиворечивости, актуальности по отношению к~решаемым 


целевым аналитическим задачам;


\item качества входной информации, подлежащей вводу: исходная текстовая 


информация, подлежащая вводу, представлена в~электронном или печатном, 


или рукописном виде, или ее нужно получить от ка\-ко\-го-то иного лица 


и~т.\,д.;


\item объема и~качества мультимедийной информации, которую пользователь 


должен обработать,


\item объема выходных данных: результатов выполнения аналитических 


расчетов, требующих от оператора последующих  


ана\-ли\-за/ин\-тер\-пре\-та\-ции/при\-ня\-тия решений.


\end{itemize}


      


      Временн$\acute{\mbox{ы}}$е затраты на выполнение автоматических операций зависят:


      \begin{itemize}


\item от реальной пропускной способности используемых в~ИАС каналов 


связи;


\item реальных показателей производительности аппаратно-программной 


платформы по при\-ему/пе\-ре\-да\-че, чте\-нию/за\-пи\-си и~хранению 


информации;


\item сложности аналитических расчетов, выполняемых в~автоматическом 


режиме;


\item объема и~структуры отчетов, формируемых в~автоматическом режиме;


\item особенностей специализированных средств, применяемых в~ИАС, 


например средств обеспечения информационной безопасности~--- 


программных средств аутентификации, средств шифрования сетевого 


трафика, антивирусных программ и~пр.


\end{itemize}


      \item Определение характеристик временн$\acute{\mbox{ы}}$х затрат на выполнение 


отдельных операций и~решение целевых задач предлагается проводить в~два 


этапа. На первом этапе рассчитывается теоретическое прогнозное значение 


соответствующей характеристики согласно предлагаемой факторной модели. 


На втором этапе в~процессе проведения нагрузочного тестирования на стенде 


или в~рамках опытной эксплуатации создаваемой ИАС производится ряд 


специализированных метрологических экспериментов по выполнению  


опе\-ра\-ций/ре\-ше\-нию задач, по результатам которых выполняется 


статистическое уточнение факторов модели. При этом следует отметить, что 


прогно\-зи\-ру\-емые значения характеристик и~их уточненные версии являются 


некоторыми опорными значениями, которые могут отличаться от их 


реальных аналогов как в~б$\acute{\mbox{о}}$льшую, так и~в меньшую сторону. Поэтому 


в~процессе промышленной эксплуатации также должна накапливаться 


статистическая информация для последующего уточнения этих 


характеристик.


      \item Ряд статистических данных, соответствующих временн$\acute{\mbox{ы}}$м 


затратам на выполнение отдельных операций, может регистрироваться 


в~системных журналах, формируемых элементами СПО ИАС. Остальные~--- 


регистрируются пользователями в~таблицах, соответствующих каждому виду 


затрат. При этом операции должны выполняться именно теми 


пользователями, которые впоследствии будут постоянно эксплуатировать 


систему. Такой подход позволит учесть уровень квалификации 


и~индивидуальные особенности обслуживающего персонала ИАС.


      \item Методика оценки временн$\acute{\mbox{ы}}$х затрат базируется:


      \begin{itemize}


\item на требованиях к~качеству и~надежности программных средств, 


регламентированных в~стандартах~[5, 6],


\item методике расчета показателей надежности, предложенной  


в~\cite{1-bor, 3-bor},


\item сведениях из математической статистики~[4, 7] для построения 


процедур уточнения характеристик трудоемкости по собранным 


статистическим данным.


\end{itemize}


\end{enumerate}





\subsection{Оценивание временных затрат на~выполнение 


неавтоматических операций}





      \subsubsection{Операции ввода текстовых данных}


      


      К этому виду относятся все операции по вво\-ду\!/\!вы\-бо\-ру из 


предопределенных списков текстовых данных в~поля различных форм.


      


      Формула вычисления трудоемкости подобной операции имеет вид:


      \begin{equation}


      T_1=\tau_1 n_1\,,


      \label{e1-bor}


      \end{equation}


где $\tau_1$~--- среднее время заполнения одного поля экранной формы; 


$n_1$~--- число заполняемых полей формы. Например, если принять 


$\tau_1\hm=15$~с, то время заполнения формы, содержащей~16~полей, 


составит~4~мин, 35~полей~--- 8~мин 45~с.


      


      Далее временн$\acute{\mbox{ы}}$е затраты на ввод текстовых данных должны быть 


скорректированы на основании статистической информации. Для получения 


этой информации привлекаются пользователи, осуществляющие предметную 


деятельность в~ИАС. Численность привлеченных сотрудников может 


составить порядка~10~человек, что является достаточным для обеспечения 


репрезентативности выборки. Для получения статистической информации 


должны быть подготовлены текстовые данные для ввода: например, 


по~10~вариантов полного набора данных для каждой из имеющихся 


экранных форм. Статистическая информация получается и~обрабатывается 


путем проведения следующей последовательности операций:


      \begin{itemize}


\item каждый испытуемый без перерыва заполняет полный набор данных 


($M N$, где $M$~--- число форм; $N$~--- число подготовленных вариантов 


данных) по предоставленным ему печатным данным. Время заполнения 


каждого набора фиксируется. Затем происходит сверка набранных данных 


и~исходной текстовой информации на предмет выявления ошибок. Число 


ошибок также фиксируется;


\item испытуемые разбиваются на пары, и~1-й номер в~паре заполняет 


полный набор данных по информации, которую ему диктует~2-й номер. 


Затем номера меняются местами и~выполняют ту же проверку. Время 


заполнения каждого набора также фиксируется. После этого происходит 


сверка набранных данных и~исходной текстовой информации на предмет 


выявления ошибок, и~число обнаруженных ошибок также фиксируется;


\item уточнение~$\hat{\tau}_1$ параметра~$\tau_1$ производится по 


следующей формуле:


\begin{equation*}


\!\hat{\tau}_1=\left(\! 1+\fr{e_\Sigma}{2NK\sum\nolimits_{m=1}^M F_m}\right)\!


\fr{1}{2NK\sum\nolimits_{m=1}^M F_m} \sum\limits^M_{m=1} 


\sum\limits^N_{n=1} \sum\limits^K_{k=1} \left( 


T^\prime_{mnk}+T_{mnk}^{\prime\prime}\right),


%\label{e2-bor}


\end{equation*}


где $M$~--- число экранных форм; $F_m$~--- число полей $m$-й формы; 


$N$~--- число подготовленных вариантов данных; $K$~--- число 


привлеченных пользователей; $e_\Sigma$~--- общее число ошибок ввода; 


$T^\prime_{mnk}$~--- время заполнения $n$-го варианта $m$-й формы $k$-м 


сотрудником по машинописному образцу; $T^{\prime\prime}_{mnk}$~--- 


время заполнения $n$-го варианта $m$-й формы $k$-м пользователем под 


диктовку.


\end{itemize}


      


      \subsubsection{Операции поиска\!/\!отбора результатов}


      


      Будем считать, что процесс поиска представляет собой 


последовательность выполнения различного рода запросов к~одной или 


нескольким таблицам баз данных ИАС. Основную часть составляют запросы 


типа


 \begin{multline*}


      \mathrm{SELECT\ \ TOP}\ \ [N]\ \ \mathrm{FROM}\ \ [\mathrm{TABLE}]\ \ \mathrm{WHERE}\  \


[\mathrm{PAR}_1=X_1]\ \ \mathrm{AND}\cdots\\


\cdots\ 


      \mathrm{AND}\ \ [\mathrm{PAR}_M=X_M]\,,


\end{multline*}


т.\,е.\ запросы на выборку~$N$ первых ответов, одновременно 


удовлетворяющих условиям $\mathrm{PAR}_i\hm= X_i$, $i\hm=1,2,\ldots, M$. Общие 


затраты на эту операцию складываются из затрат на поиск по таблицам баз 


данных, а также на формирование результата.





      Формула для опорной оценки времени имеет вид:


      \begin{equation*}


      T_2=\tau_2 MQ+\tau_3N\,,


%      \label{e3-bor}


      \end{equation*}


где $\tau_2$~--- время сравнения одного поля условия с~соответствующим 


полем в~одной строке таблицы; $M$~--- число условий в~запросе; $Q$~--- 


общее число строк в~таблицах, по которым производится поиск; $\tau_3$~--- 


время формирования одной строки ответа.





      Опорную оценку рассчитать можно, например, из таких 


предположений: $\tau_2\hm=0{,}0001$~с (с~учетом возможного выполнения 


операции чтения страниц базы данных с~диска в~оперативную память); 


$M\hm=5$ (число условий в~запросе); $Q \hm= 200\,000$ (верхняя оценка 


общего числа записей в~базе данных); $\tau_3\hm=0{,}0001$~с (время 


формирования строки ответа сравнимо с~длительностью операции  


чте\-ния\!/\!за\-пи\-си); $N \hm= 100$ (максимальная длина ответа, приемлемая 


для дальнейшего анализа). При выборе таких параметров величина $T_2\hm=  


100{,}01$~с.


      


      Показатели временн$\acute{\mbox{ы}}$х затрат на выполнение операций поиска могут 


быть значительно уточнены по результатам обработки статистики, 


накапливаемой в~процессе опытной и~промышленной эксплуатации. 


Уточнению подлежат показатели~$\tau_2$ и~$\tau_3$. Статистическая 


информация получается и~обрабатывается путем выполнения следующих 


операций:


      \begin{itemize}


\item выполняются поисковые запросы, при этом производятся измерения 


временн$\acute{\mbox{ы}}$х затрат и~фиксируются для каждого запроса вместе 


с~соответствующими значениями~$M$, $Q$, $N$ и~$T$. Суммарно должно 


быть выполнено (число~$K$) не менее~1000~запросов;


\item уточнения $\hat{\tau}_2$ и~$\hat{\tau}_3$ параметров~$\tau_2$ 


и~$\tau_3$ выполняются по формулам факторного анализа~[4]:


\begin{gather*}


\hat{\tau}_2=\fr{\overline{xz}\, \overline{y^2}-


\overline{xy}\,\overline{yz}}{\overline{x^2}\,\overline{y^2}-


(\overline{xy})^2}\,;\quad


\hat{\tau}_3=\fr{\overline{yz}\, \overline{x^2}-


\overline{xy}\,\overline{xz}}{\overline{x^2}\,\overline{y^2}-


(\overline{xy})^2}\,;\\


\displaystyle\overline{x^2}=\sum\limits^K_{k=1} x_k^2\,;\quad 


\overline{y^2}=\sum\limits^K_{k=1} y_k^2\,;\\


\displaystyle\overline{xy}=\sum\limits_{k=1}^K x_k y_k\,;\quad 


\overline{xz}=\sum\limits^K_{k=1} x_k z_k\,;\quad 


\overline{yz}=\sum\limits^K_{k=1} y_k z_k\,;\\


x_k=M_kQ_k\,;\quad y=N_k\,;\quad z_k=T_k\,.


%\label{e4-bor}


\end{gather*}


      \end{itemize}


      


      \subsubsection{Операции обработки мультимедийных данных}


      


      В качестве примеров операций по обработке мультимедиа можно 


привести выделение в~видео- или фотоматериалах изображений отдельных 


объектов, например людей, животных, зданий и~пр. Для простоты 


и~иллюстрации методики будем говорить о понятной и~распространенной 


процедуре выделения лиц людей из разных медийных источников: цифровых 


фотографий разного качества или видеопоследовательностей. 


В~подавляющем большинстве случаев лица людей в~мультимедиа могут 


быть выделены автоматически: на фото~--- практически мгновенно; на 


видео~--- в~режиме реального времени. Основные временн$\acute{\mbox{ы}}$е затраты 


связаны с~последующим отбором интересующего изображения. Модели 


временн$\acute{\mbox{ы}}$х затрат на обработку фото- и~видеоинформации при этом 


различны и~представлены ниже.


      


      Если считать, что общее среднее время на получение одного 


изображения лица прямо пропорционально числу лиц на фото, то затраты на 


эту операцию описываются следующей формулой:


      \begin{equation}


      T_3=\tau_4+n_2\tau_5\,,


      \label{e5-bor}


      \end{equation}


где $\tau_4$~--- среднее время автоматического выделения отдельных лиц 


и~сохранения выбранных фрагментов фото; $n_2$~--- среднее число лиц на 


фотографии; $\tau_5$~--- среднее время операции сравнения изображений 


(время принятия решения экспертом по одной фотографии <<да--нет>>). 


Например, при $\tau_4\hm = 20$~с, $n_2\hm = 10$ и~$\tau_5\hm = 5$~с 


оценка параметра $T_3\hm = 70$~с.





      При автоматической обработке видеопоследовательностей, как 


правило, результат выдается в~виде треков~--- последовательностей 


изображений лиц, предположительно принадлежащих одному человеку.  


Поль\-зо\-ва\-тель-экс\-перт в~этом случае дополнительно изучает 


и~изображение, и~предложенные треки, поэтому в~случае 


видеопоследовательности соотношение~(\ref{e5-bor}) уточняется для учета 


продолжительности видео и~характеристик треков:


      \begin{equation}


      T_4=\tau_6+S_1+\tau_8 n_5\,,\quad S_1=n_3 n_4 \tau_7\,,


      \label{e6-bor}


      \end{equation}


где $\tau_6$~--- длительность видеопоследовательности; $n_3$~--- число 


кадров, содержащих изображения лиц; $n_4$~--- среднее число лиц, 


встречающихся в~кадре видеопоследовательности; $\tau_7$~--- среднее время 


операции сравнения изображений; $n_5$~--- среднее число изображений 


лица в~одном треке; $\tau_8$~--- среднее время операции сравнения 


фотографий трека. Например, при $\tau_6\hm = 600$~с, $n_3\hm = 600$, 


$n_4\hm = 10$, $\tau_7\hm = 20$~с, $n_5\hm = 10$ и~$\tau_8\hm = 5$~с 


оценка параметра $T_4\hm = 1850$~с.





      Параметры, используемые в~формулах~(\ref{e5-bor})  


или~(\ref{e6-bor}), должны уточняться в~процессе эксплуатации. Процедура 


уточнения заключается в~выполнении следующей последовательности 


операций:


      \begin{itemize}


\item не менее $K\hm = 100$~раз выполняется обработка фотографий 


и~фиксируются замеры~$T_3^k$, $\tau_4^k$ и~$n_2^k$ для каждого $k$-го 


фото;


\item также не менее $K\hm = 100$~раз для видео разной продолжительности 


выполняется обработка последовательностей и~фиксируются замеры 


$T_4^k$, $\tau_6^k$, $n_3^k$, $n_4^k$, $n_5^k$ и~$S_1^k$ для каждого $k$-го 


видео;


\item параметры $n_2$, $\tau_4$ и~$\tau_5$ формулы~(\ref{e5-bor}) уточняются 


величинами~$\hat{n}_2$, $\hat{\tau}_4$ и~$\hat{\tau}_5$, 


а~параметры~$\tau_6$, $\tau_7$ и~$\tau_8$ формулы~(\ref{e6-bor}) уточняются 


величинами~$\hat{\tau}_6$, $\hat{\tau}_7$ и~$\hat{\tau}_8$ следующим 


образом:





\noindent


\begin{alignat*}{3}


\displaystyle\hat{n}_2&=\fr{1}{K}\sum\limits^K_{k=1} n_2^k\,;&\quad


\hat{\tau}_4&=\fr{1}{K}\sum\limits^K_{k=1} \tau_4^k\,;&\quad


\hat{\tau}_5&=\fr{1}{K}\sum\limits^K_{k=1} \fr{T_3^k-\tau_4^k}{n_2^k}\,;\\


\displaystyle\hat{\tau}_6&=\fr{1}{K} \sum\limits^K_{k=1} \tau_6^k\,;&\quad


\hat{\tau}_7&=\fr{1}{K}\sum\limits^K_{k=1} \fr{S_1^k}{n_3^k n_4^k}\,;&\quad


\hat{\tau}_8&=\fr{1}{K}\sum\limits^K_{k=1} \fr{T_4^k -\tau_6^k -


S_1^k}{n_5^k}\,.


%\label{e7-bor}


\end{alignat*}


      \end{itemize}


      


      \subsubsection{Операции анализа\!/\!интерпретации\!/\!принятия решений}


      


      Временн$\acute{\mbox{ы}}$е затраты на выполнение таких операций складываются из 


фиксированных затрат на принятие решений и~переменных затрат на 


выполнение операций анализа и~интерпретации данных. Переменный 


характер этих затрат связан с~вариацией объема данных, которые должны 


быть проанализированы. Уместно считать, что общее время, затрачиваемое 


на выполнение таких операций, прямо пропорционально объему отобранных 


результатов, поэтому затраты предлагается описывать следующей формулой:


      \begin{equation}


      T_5=\tau_9+n_6\tau_{10}\,,


      \label{e8-bor}


      \end{equation}


где $\tau_9$~--- среднее время, затрачиваемое на процедуру принятия 


решения; $n_6$~--- число отобранных результатов поиска; $\tau_{10}$~--- 


среднее время выполнения операции анализа и~интерпретации одного 


результата в~отобранном списке. Например, при $\tau_9\hm = 2$~мин, 


$n_6\hm = 10$ и~$\tau_{10}\hm = 0{,}5$~мин оценка параметра $T_5\hm = 


7$~мин.





      Как и~в предыдущих случаях, параметры~$\tau_9$ и~$\tau_{10}$ 


в~(\ref{e8-bor}) должны уточняться величинами~$\hat{\tau}_9$ 


и~$\hat{\tau}_{10}$: накапливается статистическая информация~$T_9^k$ 


и~$n_6^k$ в~не менее чем $K\hm = 100$ экспериментах и~факторы модели 


корректируются по следующим формулам:


      \begin{alignat*}{3}


            \displaystyle\hat{\tau}_{10}&= \fr{\overline{n_6T_5}-\overline{n_6} 


\overline{T}_5}{\overline{n_6^2} -\overline{n_6}^2}\,;&\quad


      \hat{\tau}_9 &=\overline{T}_5 -\hat{\tau}_{10} \overline{n_6}\,;&\quad


      \overline{n_6T_5}&=\fr{1}{K} \sum\limits^K_{k=1} n_6^k T_5^k\,;\\


          \displaystyle  \overline{n_6} &=\fr{1}{K}\sum\limits^K_{k=1} n_6^k\,;&\quad


      \overline{T_5}&=\fr{1}{K}\sum\limits^K_{k=1} T_5^k\,;&\quad


      \overline{n_6^2}&=\fr{1}{K}\sum\limits^K_{k=1} \left( n_6^k\right)^2\,.


      %      \label{e9-bor}


      \end{alignat*}


      


\subsection{Оценивание временных затрат на~выполнение автоматических 


операций}





      \subsubsection{Передача данных по~телекоммуникационной сети}


      


      Временн$\acute{\mbox{ы}}$е затраты на выполнение этой операции вычисляются по 


формуле:


      \begin{equation*}


      T_6=r_1 v_1\,,


      %\label{e10-bor}


      \end{equation*}


где $r_1$~--- фактическая скорость передачи данных по 


телекоммуникационным каналам; $v_1$~--- объем передаваемых данных. 


Время, затрачиваемое на выполнение этой операции, является 


незначительным по сравнению со временем, затрачиваемым на другие 


операции. Статистическое уточнение характеристик для этой операции,  


по-ви\-ди\-мо\-му, не требуется.





\vspace*{-3pt}


      


      \subsubsection{Логистические операции}


      


      К логистическим относятся такие операции, как  


за\-груз\-ка\!/\!вы\-груз\-ка данных в\!/\!из различных баз данных в~составе ИАС, 


формирование файлов отчетов, им\-порт\!/\!экс\-порт данных внешних 


информационных источников и~т.\,п. Содержательно таких операций может 


быть множество, но все они объединены общим технологическим смыслом, 


поэтому и~оценка времени для них унифицирована:





\vspace*{2pt}





\noindent


      \begin{equation}


      T_7^i= \tau^i_{11} n_7^i+\tau^i_{12}\,,


      \label{e11-bor}


      \end{equation}


      


      \vspace*{-2pt}


      


      \noindent


где $T_7^i$~--- время на выполнение логистической операции $i$-го типа; 


$\tau^i_{11}$~--- время выполнения логистической операции $i$-го типа для 


одной записи; $n_7^i$~--- объем <<пакета заданий>>; $\tau^i_{12}$~--- 


базовые затраты на выполнение логистической операции $i$-го типа. 


Например, в~предположении, что $\tau_{11}^i\hm= 0{,}5$~с, $n_7^i\hm = 


100$ и~$\tau^i_{12}\hm = 10$~с, оценка параметра $T_7^i\hm = 60$~с.





      Далее параметры $\tau^i_{11}$ и~$\tau^i_{12}$ из~(\ref{e11-bor}) 


должны уточняться величинами~$\hat{\tau}^i_{11}$ и~$\hat{\tau}^i_{12}$. 


Накапливается статистическая информация $n_7^{i,k}$ и~$T_7^{i,k}$ по 


выборке объема не менее, чем $K_i\hm=100$ для операции каждого $i$-го 


типа, а~затем факторы модели корректируются согласно следующим 


формулам:





\noindent


      \begin{alignat*}{2}


                \displaystyle  


          \hat{\tau}^i_{11}&=\fr{\overline{n_7^i T_7^i} -\overline{n_7^i}\, 


\overline{T_7^i}}{\overline{(n_7^i)^2}-\overline{n_7^i}^2}\,;&\quad


\hat{\tau}^i_{12}&=\overline{T_7^i} -\hat{\tau}^i_{11} 


\overline{n_7^i}\,;\\


      \overline{n_7^i T_7^i}&=\fr{1}{K_i}\sum\limits_{k=1}^{K_i} n_7^{i,k} 


T_7^{i,k}\,;&\quad \overline{n_7^i}&=\fr{1}{K}\sum\limits_{k=1}^{K_i} 


n_7^{i,k}\,;\\


\displaystyle  \overline{T_7^i} &=\fr{1}{K_i} \sum\limits_{k=1}^{K_i} 


T_7^{i,k}\,;&\quad \overline{(n_7^i)^2} &=\fr{1}{K_i} \sum\limits_{k=1}^{K_i} 


\left( n_7^{i,k}\right)^2\,.


      %      \label{e12-bor}


      \end{alignat*}


      


      \vspace*{-6pt}


      


      \subsubsection{Аналитические расчеты и~формирование отчетов}


      


      Оставшиеся два типа автоматических операций, по мнению авторов 


статьи, в~каждой ИАС носят слишком специфический характер, чтобы давать 


по ним ка\-кие-ли\-бо универсальные рекомендации. Конкретные алгоритмы как 


при проведении аналитических расчетов, так и~при формировании отчетов 


могут оказаться и~довольно сложными вычислительными процессами, 


и~простейшими процедурами сохранения результатов поиска в~файл. 


Наиболее уместным в~этом случае представляется не проводить  


ка\-ких-ли\-бо предварительных расчетов, а~просто набрать достаточную 


статистику в~процессе ис\-пы\-та\-ний/экс\-плу\-а\-та\-ции и~применить 


усредненные оценки или вообще ограничиться экспертным мнением 


о~данных параметрах.


      


\section{Методика оценивания численности персонала}





\subsection{Общие принципы оценивания численности}





      В предыдущем разделе были оценены временн$\acute{\mbox{ы}}$е характеристики 


выполнения отдельных операций и~решения задач в~типовой ИАС. 


Полученные результаты являются основой вычисления следующих 


характеристик численности обслуживающего персонала, необходимого для 


решения некоторого фиксированного объема задач.


      


      \textit{Производительность}~$r$ [операций\!/\!ед.\ вр.] 


\textit{пользователя по выполнению од\-ной операции одного определенного 


типа} вычисляется по формуле:


      \begin{equation}


      r=\fr{1}{\tau}\,,


      \label{e13-bor}


      \end{equation}


где $\tau$ [ед.\ вр.]~--- среднее время выполнения одной операции (например, 


$T_1,\ldots , T_5$ из предыдущего раздела).


      


      \textit{Время $T_{\mathrm{nec}}$, необходимое для выполнения~$N$~операций 


одного определенного типа~$n$ операторами}, вычисляется по формуле:


      \begin{equation}


      T_{\mathrm{nec}}= \left\lceil \fr{N}{n}\right\rceil\,\tau\,,


      \label{e14-bor}


      \end{equation}


где $\lceil\cdot\rceil$~--- операция округления в~б$\acute{\mbox{о}}$льшую сторону.





      \textit{Число операторов, необходимое для выполнения~$N$~операций 


за время~$T$}, определяется формулой:


      \begin{equation}


      n_{\mathrm{nec}}= \begin{cases}


      \mbox{не\ определено}, & \mbox{если } T<\tau\,;\\


      \left\lceil \fr{N\tau}{T}\right\rceil\,, & \mbox{если } T\geq\tau\,.


      \end{cases}


      \label{e15-bor}


      \end{equation}


      


\subsection{Пример оценки численности персонала}





      Для оценки численности пользовательского персонала конкретной 


ИАС необходимо формально определить перечень функций, которые будут 


исполняться персоналом. Этот перечень должен быть согласован с~расчетами 


временн$\acute{\mbox{ы}}$х затрат, выполненных в~предыдущей части работы, т.\,е.\ либо 


основываться на той же терминологии, либо на близком аналоге или 


комбинации рас\-смот\-рен\-ных понятий. Далее приведен простой пример


 одного вида деятельности,  поясняющий предлагаемую методику. Для других 


видов деятельности расчеты могут быть выполнены аналогично.


%\linebreak\vspace*{-12pt}





\pagebreak





\





\vspace*{-12pt}





\renewcommand{\figurename}{\protect\bf Таблица}





\begin{floatingfigure}[t]{64mm}


\vspace*{-18pt}


{\small


\begin{center}


\Caption{Модельные расчеты времени~$T_{\mathrm{nec}}$}


\vspace{2ex}





\begin{tabular}{|c|c|c|c|}


\hline


&\multicolumn{3}{c|}{\ }\\[-9pt]


\multicolumn{1}{|c|}{\raisebox{-6pt}[0pt][0pt]{$n$}}&


\multicolumn{3}{c|}{$T_{\mathrm{nec}}$, ч}\\


%&&&\\[-9pt]


\cline{2-4}


&&&\\[-9pt]


&$N = 10$&$N = 100$&$N = 1000$\\


\hline


1&0,8\hphantom{9}&8\hphantom{,9}&80\\


2&0,4\hphantom{9}&4\hphantom{,9}&40\\


5&0,16&1,6&16\\


10\hphantom{9}&0,08&0,8&\hphantom{9}8\\


\hline


\end{tabular}


\end{center}}


%\end{table}


{\small


\begin{center}


\Caption{Модельные расчеты числа операторов~$n_{\mathrm{nec}}$}


\vspace*{2ex}





\begin{tabular}{|c|c|c|c|}


\hline


\multicolumn{1}{|c|}{\raisebox{-6pt}[0pt][0pt]{$T$, ч}}&\multicolumn{3}{c|}


{$n_{\mathrm{nec}}$, ч}\\


%&&&\\[-9pt]


\cline{2-4}


&&&\\[-9pt]


&$N = 10$&$N = 100$&$N = 1000$\\


\hline


\hphantom{9}1&1&8&80\\


\hphantom{9}8&1&1&10\\


24&1&1&\hphantom{9}2\\


40&1&1&\hphantom{9}1\\


\hline


\end{tabular}


\end{center}}


\vspace*{6pt}


\end{floatingfigure}

















      


      В процессе эксплуатации любой ИАС выполняются операции 


администрирования, одной из которых является регистрация нового 


пользователя системы. Она относится к~типу операций ввода атрибутивной 


информации, и~ее трудоемкость вычисляется по формуле~(\ref{e1-bor}). 


Число заполняемых полей в~этой операции может варьироваться 


в~зависимости от нюансов разграничения доступа в~конкретной ИАС. Если 


принять $\tau_1\hm =15$~с и~$n_1\hm=20$, то время регистрации одного 


пользователя составит~5~мин, а~начальная оценка~(\ref{e13-bor}) 


производительности~$r$~--- 12~операций/ч.


      


      В табл.~1 и~2 приведены результаты выполненных модельных расчетов 


количества времени~$T_{\mathrm{nec}}$, необходимого для 


регистрации~$N$~пользователей $n$~администраторами, вычисленные по 


формуле~(\ref{e14-bor}), и~модельных расчетов числа операторов~$n_{\mathrm{nec}}$, 


необходимых для регистрации~$N$~пользователей за время~$T$, 


вычисленные по формуле~(\ref{e15-bor}).











      








\renewcommand{\figurename}{\protect\bf Рис.}


\renewcommand{\tablename}{\protect\bf Таблица}





\vspace*{-4pt}





\section{Заключение}





\vspace*{-2pt}





      В данной статье представлена методика предварительного расчета 


и~статистического уточнения параметров трудоемкости решения целевых 


задач ИАС и~численности пользовательского персонала. Ее ключевые 


преимущества~--- операционное описание решения целевых задач ИАС 


и~простота используемых математических моделей, а также возможность их 


корректировки по данным эксплуатации конкретной ИАС.


      


      Недостатки представленной методики являются в~некотором смысле 


продолжением ее достоинств. Универсальность идеи декомпозиции 


подразумевает детальную схему <<сборки>> алгоритма решения конкретной 


целевой задачи из набора имеющихся операций. Поэтому повышение 


точности оценок характеристик временн$\acute{\mbox{ы}}$х затрат на выполнение функций 


назначения ИАС связано с~уточнением описания этих целевых функций для 


конкретных ИАС. Одновременно с~этим математические модели определения 


временн$\acute{\mbox{ы}}$х затрат на выполнение отдельных операций могут быть 


значительно уточнены за счет включения в~них других факторов. Оба эти 


направления определяют возможные перспективы дальнейших 


исследований.


      


{\small\frenchspacing


{%\baselineskip=10.8pt


\addcontentsline{toc}{section}{Литература}


\begin{thebibliography}{9}





\bibitem{2-bor}


\Au{Chamoni P., Gluchowski~P.} Analytische Informationssysteme: Business 


Intelligence-Technologien und~--- Anwendungen.~--- Berlin: Springer, 2004. 470~p.


\bibitem{3-bor}


\Au{Борисов А.\,В., Босов А.\,В., Иванов~А.\,В., Корепанов~Э.\,Р.} К~вопросу расчета 


надежности ин\-фор\-ма\-ци\-он\-но-те\-ле\-ком\-му\-ни\-ка\-ци\-он\-ных систем: учет 


характеристик программного обеспечения~/\!\!/ Системы и~средства информатики, 2018. 


Т.~28. №\,1. С.~20--34.





\bibitem{1-bor} %3


\Au{Борисов А.\,В., Босов~А.\,В., Иванов~А.\,В., Чавтараев~Р.\,Б.} Имитационное 


моделирование пользовательской активности для оценивания вероятностно-временных 


характеристик программного обеспечения~/\!\!/ Системы и~средства информатики, 2018. 


Т.~28. №\,2. С.~20--33.


\bibitem{4-bor}


\Au{Афифи А., Эйзен С.} Статистический анализ: Подход с~использованием ЭВМ~/


Пер с~англ.~--- М.: 


Мир, 1982. 488~c. (\Au{Afifi~A.\,A., Azen~S.\,P.} {Statistical analysis: 


A~computer oriented approach}.~--- 2nd ed.~--- New York, NY,


USA: Academic Press, 1979. 442~p.)


\bibitem{5-bor}


ГОСТ 28806-90. Межгосударственный стандарт. Качество программных средств.~--- М.: 


Изд-во стандартов, 2001. 8~с.


\bibitem{6-bor}


ГОСТ Р ИСО/МЭК 9126-93. Информационная технология. Оценка программной 


продукции характеристики качества и~руководства по их применению.~--- М.: Изд-во 


стандартов, 2004. 10~с.


\bibitem{7-bor}


\Au{Ширяев А.\,Н.} Вероятность.~--- М.: Наука, 1979. 574~с.





        \end{thebibliography}


} }








\vspace*{-6pt}





\hfill{\small\textit{Поступила в~редакцию 19.01.18}}














\vspace*{8pt}





\hrule





\vspace*{2pt}





%\newpage





%\vspace*{-28pt}





\hrule





%\vspace*{-3pt}











\def\tit{ANALYTICAL INFORMATION SYSTEMS PERFORMANCE 


ANALYSIS: METHODOLOGY FOR~TIMETABLE AND~STAFF~QUANTITY~EVALUATION}











\def\titkol{Analytical information systems performance 


analysis: Methodology for timetable evaluation} %


% and~staff  quantity evaluation}








\def\aut{A.\,V.~Borisov, A.\,V.~Bosov, A.\,V.~Ivanov, 


and~A.\,I.~Stefanovich}


\def\autkol{A.\,V.~Borisov, A.\,V.~Bosov, A.\,V.~Ivanov, 


and~A.\,I.~Stefanovich}











\titel{\tit}{\aut}{\autkol}{\titkol}





\vspace*{-9pt}








 \noindent





\def\leftfootline{\small{\textbf{\thepage}


\hfill Sistemy i Sredstva Informatiki~---


Systems and Means of Informatics 2018 vol~28 no\,3}


}%


 \def\rightfootline{\small{Sistemy i Sredstva Informatiki~---


 Systems and Means of Informatics   2018   vol~28   no\,3


\hfill \textbf{\thepage}}}








     


      \noindent


      Institute of Informatics Problems, Federal Research Center ``Computer 


Science and~Control'' of the Russian Academy of Sciences, 44-2~Vavilov Str., 


Moscow 119333, Russian Federation





\vspace*{-3pt}


        


        \Abste{The approach  proposed earlier to evaluate the 


        software performance based on  simulation and statistical processing 


        of experimental results is developed for the needs of staff management 


        of an analytical information system (AIS). To this purpose, staff 


        responsibilities are classified to decompose the major AIS


        tasks 


        solved by the corresponding application software. The paper presents 


        methodical recommendations how to choose a~set of timetable 


        characteristics and observable metrics and collect related statistical 


        information during the AIS testing and operation. The paper also contains 


        the algorithms for calculation of the parameters of operations timetable 


        and their statistical refinement. The


        methodology for assessing the 


        performance of staff and the 


        required number of staff is illustrated by simple examples.}


        


        \KWE{analytical information system; probability characteristics; 


        statistical methods; factor analysis}








\DOI{10.14357\!/\!08696527180303} 





%\vspace*{-24pt}





%\Ack


%\noindent





 \vspace*{-12pt}








\renewcommand{\bibname}{\protect\rmfamily References}


%\renewcommand{\bibname}{\large\protect\rm References}





{\small\frenchspacing


{%\baselineskip=10.8pt


\addcontentsline{toc}{section}{References}


\begin{thebibliography}{9}





\bibitem{2-bor-1}


\Aue{Chamoni, P., and P.~Gluchowski}. 2004. \textit{Analytische 


Informationsysteme: Business Intelligence-Technologien und~--- Anwendungen}. 


Berlin: Springer. 470~p.


\bibitem{3-bor-1}


\Aue{Borisov, A.\,V., A.\,V.~Bosov, A.\,V.~Ivanov, and E.\,R.~Korepanov.} 


2018. K~voprosu rascheta nadezhnosti informatsionno-telekommunikatsionnykh 


sistem: uchet kha\-rak\-te\-ristik programmnogo obespecheniya [To the reliability of 


information-telecommunication system: An approach to recognition of reliable 


software characteristics]. \textit{Sistemy i~Sredstva Informatiki~--- Systems and 


Means of Informatics} 28(1):20--34.





\bibitem{1-bor-1}


\Aue{Borisov, A.\,V., A.\,V.~Bosov, A.\,V.~Ivanov, and R.\,V.~Chavtaraev}. 


2018. Imitatsionnoe modelirovanie pol'zovatel'skoy aktivnosti dlya otsenivaniya 


veroyatnostno-vremennykh kharakteristik programmnogo obespecheniya [Monte 


Carlo based user activity simulation for software performance evaluation]. 


\textit{Sistemy i~Sredstva Informatiki~--- Systems and Means of Informatics} 


28(2):20--33.


\bibitem{4-bor-1}


\Aue{Afifi, A.\,A., and S.\,P.~Azen.} 1979. \textit{Statistical analysis: 


A~computer oriented approach}. 2nd ed. New York, NY: Academic Press. 442~p.


\bibitem{5-bor-1}


GOST 28806-90. 2001. Mezhgosudarstvennyy standart. Kachestvo programmnykh 


sredstv [Software quality. Terms and definitions]. Moscow: Standards Publs. 8~p. 


\bibitem{6-bor-1}


GOST R ISO/IEC 9126-93. 2004. Informatsionnaya tekhnologiya. Otsenka 


programmnoy produktsii kharakteristiki kachestva i rukovodstva po ikh 


primeneniyu [Information technology. Software product evaluation. Quality 


characteristics and guidelines for their use]. Moscow: Standards Publs.


10~p.


\bibitem{7-bor-1}


\Aue{Shiryaev, A.\,N.} 1984. \textit{Probability}. New York, NY: Springer-Verlag. 


578~p.








\end{thebibliography}


} }








\vspace*{-4pt}





\hfill{\small\textit{Received January 19, 2018}}    





\vspace*{-9pt}


  





\Contr





\vspace*{-4pt}





\noindent


\textbf{Borisov Andrey V.} (b.\ 1965)~--- Doctor of Science in physics and 


mathematics, principal scientist, Institute of Informatics Problems, Federal 


Research Center ``Computer Science and Control'' of the Russian Academy 


of Sciences, 44-2~Vavilov Str., Moscow 119333, Russian Federation; 


\mbox{ABorosov@ipiran.ru}





\vspace*{2pt}





\noindent


\textbf{Bosov Alexey V.} (b.\ 1969)~--- Doctor of Science in technology, principal scientist, Institute of 


Informatics Problems, Federal Research Center ``Computer Science and Control'' of the Russian 


Academy of Sciences, 44-2~Vavilov Str., Moscow 119333, Russian Federation; 


\mbox{AVBosov@ipiran.ru}





\vspace*{2pt}





\noindent


\textbf{Ivanov Alexey V.} (b.\ 1976)~--- scientist, Institute of Informatics Problems, Federal Research 


Center ``Computer Science and Control'' of the Russian Academy of Sciences, 44-2~Vavilov Str., Moscow 


119333, Russian Federation; \mbox{AIvanov@ipiran.ru}





\vspace*{2pt}





\noindent


\textbf{Stefanovich Alexey I.} (b.\ 1983)~--- junior scientist, Institute of Informatics Problems, Federal 


Research Center ``Computer Science and Control'' of the Russian Academy of Sciences, 44-2~Vavilov 


Str., Moscow 119333, Russian Federation; \mbox{AStefanovich@frccsc.ru}


\label{end\stat}








\renewcommand{\bibname}{\protect\rm Литература}


